% GENERATED FILE. Edit canonical lessons, then run npm run generate:books.
% canonical-source-hash: fnv1a64:406116a9f087cd6b
% canonical-lessons: TE-C170-samayam, TE-C170-appudappudu, TE-C170-taracuga, TE-C170-tedi, TE-C170-varantam

\chapter{Time, Sometimes, Often, Date, Weekend}
\label{ch:te-time-sometimes-often-date-weekend}
\hlchaptermodality{170}

\begin{chapteropening}
\textbf{By the end of this chapter:} \emph{I can say time, sometimes, often, a date and a weekend.}

Complete the last lesson of chapter 170: I can say time, sometimes, often, a date and a weekend.
\end{chapteropening}

\section[\texorpdfstring{\te{సమయం} (samayaṁ)}{సమయం (samayaṁ)}]{\te{సమయం} (samayaṁ) — time}

\label{lesson:TE-C170-samayam}



\begin{quote}
Before the new one: say the Telugu for a scar, then the Telugu for a sprain.
\end{quote}

\subsection*{You'll want to know: \te{సమయం}}
\textbf{\te{సమయం}} — \emph{samayaṁ} — \textquotedblleft{}time\textquotedblright{}.

\begin{grammarlens}[title={What you've built}]
One more word for this chapter.
\end{grammarlens}

\subsection*{Your turn}
\begin{itemize}
\raggedright
  \item \emph{Say it:} \emph{samayaṁ}
  \item \emph{Say it:} \emph{samayaṁ}, once more
  \item \emph{Say it:} say \emph{samayaṁ}
  \item \emph{Recall:} say the Telugu for a scar, then the Telugu for a sprain, then say \emph{samayaṁ} again
\end{itemize}

\begin{tcolorbox}[breakable,colback=teal!4,colframe=teal!35!black,title={Before you move on}]
What does \te{సమయం} mean? (\textquotedblleft{}Time\textquotedblright{}.) Say it once more.
\end{tcolorbox}
\section[\texorpdfstring{\te{అప్పుడప్పుడు} (appuḍappuḍu)}{అప్పుడప్పుడు (appuḍappuḍu)}]{\te{అప్పుడప్పుడు} (appuḍappuḍu) — sometimes}

\label{lesson:TE-C170-appudappudu}



\begin{quote}
Before the new one: say the Telugu for a sprain, then the Telugu for time.
\end{quote}

\subsection*{You'll want to know: \te{అప్పుడప్పుడు}}
\textbf{\te{అప్పుడప్పుడు}} — \emph{appuḍappuḍu} — \textquotedblleft{}sometimes\textquotedblright{}.

\begin{grammarlens}[title={What you've built}]
One more word for this chapter.
\end{grammarlens}

\subsection*{Your turn}
\begin{itemize}
\raggedright
  \item \emph{Say it:} \emph{appuḍappuḍu}
  \item \emph{Say it:} \emph{appuḍappuḍu}, once more
  \item \emph{Say it:} say \emph{appuḍappuḍu}
  \item \emph{Recall:} say the Telugu for a sprain, then the Telugu for time, then say \emph{appuḍappuḍu} again
\end{itemize}

\begin{tcolorbox}[breakable,colback=teal!4,colframe=teal!35!black,title={Before you move on}]
What does \te{అప్పుడప్పుడు} mean? (\textquotedblleft{}Sometimes\textquotedblright{}.) Say it once more.
\end{tcolorbox}
\section[\texorpdfstring{\te{తరచుగా} (taracugā)}{తరచుగా (taracugā)}]{\te{తరచుగా} (taracugā) — often}

\label{lesson:TE-C170-taracuga}



\begin{quote}
Before the new one: say the Telugu for time, then the Telugu for sometimes.
\end{quote}

\subsection*{You'll want to know: \te{తరచుగా}}
\textbf{\te{తరచుగా}} — \emph{taracugā} — \textquotedblleft{}often\textquotedblright{}.

\begin{grammarlens}[title={What you've built}]
One more word for this chapter.
\end{grammarlens}

\subsection*{Your turn}
\begin{itemize}
\raggedright
  \item \emph{Say it:} \emph{taracugā}
  \item \emph{Say it:} \emph{taracugā}, once more
  \item \emph{Say it:} say \emph{taracugā}
  \item \emph{Recall:} say the Telugu for time, then the Telugu for sometimes, then say \emph{taracugā} again
\end{itemize}

\begin{tcolorbox}[breakable,colback=teal!4,colframe=teal!35!black,title={Before you move on}]
What does \te{తరచుగా} mean? (\textquotedblleft{}Often\textquotedblright{}.) Say it once more.
\end{tcolorbox}
\section[\texorpdfstring{\te{తేదీ} (tēdī)}{తేదీ (tēdī)}]{\te{తేదీ} (tēdī) — a date}

\label{lesson:TE-C170-tedi}



\begin{quote}
Before the new one: say the Telugu for sometimes, then the Telugu for often.
\end{quote}

\subsection*{You'll want to know: \te{తేదీ}}
\textbf{\te{తేదీ}} — \emph{tēdī} — \textquotedblleft{}a date\textquotedblright{}.

\begin{grammarlens}[title={What you've built}]
One more word for this chapter.
\end{grammarlens}

\subsection*{Your turn}
\begin{itemize}
\raggedright
  \item \emph{Say it:} \emph{tēdī}
  \item \emph{Say it:} \emph{tēdī}, once more
  \item \emph{Say it:} say \emph{tēdī}
  \item \emph{Recall:} say the Telugu for sometimes, then the Telugu for often, then say \emph{tēdī} again
\end{itemize}

\begin{tcolorbox}[breakable,colback=teal!4,colframe=teal!35!black,title={Before you move on}]
What does \te{తేదీ} mean? (\textquotedblleft{}A date\textquotedblright{}.) Say it once more.
\end{tcolorbox}
\section[\texorpdfstring{\te{వారాంతం} (vārāntaṁ)}{వారాంతం (vārāntaṁ)}]{\te{వారాంతం} (vārāntaṁ) — a weekend}

\label{lesson:TE-C170-varantam}



\begin{quote}
Before the new one: say the Telugu for often, then the Telugu for a date.
\end{quote}

\subsection*{You'll want to know: \te{వారాంతం}}
\textbf{\te{వారాంతం}} — \emph{vārāntaṁ} — \textquotedblleft{}a weekend\textquotedblright{}.

\begin{grammarlens}[title={What you've built}]
One more word for this chapter.
\end{grammarlens}

\subsection*{Your turn}
\begin{itemize}
\raggedright
  \item \emph{Say it:} \emph{vārāntaṁ}
  \item \emph{Say it:} \emph{vārāntaṁ}, once more
  \item \emph{Say it:} say \emph{vārāntaṁ}
  \item \emph{Recall:} say the Telugu for often, then the Telugu for a date, then say \emph{vārāntaṁ} again
\end{itemize}

\begin{tcolorbox}[breakable,colback=teal!4,colframe=teal!35!black,title={Before you move on}]
What does \te{వారాంతం} mean? (\textquotedblleft{}A weekend\textquotedblright{}.) Say it once more.
\end{tcolorbox}
